% !TEX program = lualatex
\documentclass[professionalfonts]{beamer}
\newcommand{\slidemode}{1}  % カラー投影モード
\usepackage{beamer_template}
\usepackage{enumerate}

\newcommand{\LectureNum}{2}
\newcommand{\LectureTitle}{相似性と移動法則}
\newcommand{\CourseName}{応用数学}
\renewcommand{\TermName}{前期}

\begin{document}

\begin{frame}{第2回　相似性と移動法則}
  \begin{block}{本回の内容}
  相似性と移動法則
  \end{block}
  \vfill\hfill{\scriptsize \textcolor{gray}{第2章 ラプラス変換　§1}}
\end{frame}

\begin{frame}{定理：相似性}
  \begin{block}{}
  \[
    L[f(at)] = (1/a) F(s/a) \ (a は正の定数)
  \]
  \end{block}
\end{frame}

\begin{frame}{定理：像関数の移動法則}
  \begin{block}{}
  \[
    L[e^{at}f(t)] = F(s-a) \ (a は定数)
  \]
  \end{block}
\end{frame}

\begin{frame}{定理：原関数の移動法則（ヘビサイドの移動定理）}
  \begin{block}{}
  \[
    L[f(t-\mu)U(t-\mu)] = e^{-\mu s}F(s) \ (\mu > 0)
  \]
  \end{block}
\end{frame}

\begin{frame}{例題7}
  \begin{exampleblock}{問題}
    $L[\sin(\omega t)]$ を求めよ（ $\omega$ は 0 でない定数）\\[2pt]
  \end{exampleblock}
\end{frame}

\begin{frame}{例題7　解答}
  $\omega > 0$ のとき : 相似性と例題 4 の結果から\\[4pt]
  $L[\sin(\omega t)] = (1/\omega) \cdot 1/((s/\omega)^{2}+1) = \omega/(s^{2}+\omega^{2})$\\[4pt]
  $\omega < 0$ のとき $: \sin(-\omega t) = -\sin(\omega t)$ を利用\\[4pt]
  $L[\sin(\omega t)] = -L[\sin(-\omega t)] = -(-\omega)/(s^{2}+\omega^{2}) = \omega/(s^{2}+\omega^{2})$
  \\[8pt]\textcolor{red}{\textbf{答}}：$L[\sin(\omega t)] = \omega/(s^{2}+\omega^{2}) \ (s > 0)$（$\omega$ の正負に関係なく）
\end{frame}

\begin{frame}{例題8}
  \begin{exampleblock}{問題}
    $L[e^{\alpha t}\sin(\beta t)]$ を求めよ（ $\alpha$ は定数、 $\beta \neq 0$ ）\\[2pt]
  \end{exampleblock}
\end{frame}

\begin{frame}{例題8　解答}
  例題 7 より $L[\sin(\beta t)] = \beta/(s^{2}+\beta^{2}) \ (s > 0)$\\[4pt]
  像関数の移動法則 $L[e^{\alpha t}f(t)] = F(s-\alpha)$ を適用すると\\[4pt]
  s を $s-\alpha$ に置き換えて\\[4pt]
  $L[e^{\alpha t}\sin(\beta t)] = \beta/((s-\alpha)^{2}+\beta^{2})$
  \\[8pt]\textcolor{red}{\textbf{答}}：$L[e^{\alpha t}\sin(\beta t)] = \beta/((s-\alpha)^{2}+\beta^{2}) \ (s > \alpha)$
\end{frame}

\begin{frame}{演習問題}
  \begin{exampleblock}{自分で解いてみよう}
  \begin{description}
    \item[問8] $L[\cos(\omega t)]$ を求めよ（ $\omega$ は 0 でない定数）
    \item[問9] $\sin^{2}t, \cos^{2}t$ のラプラス変換を求めよ（半角の公式を使用）
    \item[問10] $L[te^{\alpha t}]$, $L[e^{\alpha t}\cos(\beta t)]$ を求めよ（ $\alpha, \beta$ は定数、$\beta \neq 0$）
    \item[問11] $L[\sin(t - \pi/4)\cdot U(t - \pi/4)]$ を求めよ
    \item[問12] 関数 $f(t) = (t-2)U(t-2)$ のグラフをかけ。また $L[f(t)]$ を求めよ
  \end{description}
  \end{exampleblock}
\end{frame}

\begin{frame}{問8　解答}
  相似性 $: L[f(at)] = (1/a)F(s/a), F(s) = L[\cos t] = s/(s^{2}+1)$ を利用\\[4pt]
  $L[\cos(\omega t)] = (1/\omega)\cdot(s/\omega)/((s/\omega)^{2}+1) = (1/\omega)\cdot(s/\omega)\cdot\omega^{2}/(s^{2}+\omega^{2})$
  \\[8pt]\textcolor{red}{\textbf{答}}：$L[\cos(\omega t)] = s/(s^{2}+\omega^{2}) \ (s > 0)$（$\cos$ は偶関数なので $\omega$ の正負に関係なく成立）
\end{frame}

\begin{frame}{問9　解答}
  $\sin^{2}t = (1-\cos 2t)/2   \to   L[\sin^{2}t] = (1/2)(1/s - s/(s^{2}+4)) = 2/(s(s^{2}+4))$\\[4pt]
  $\cos^{2}t = (1+\cos 2t)/2   \to   L[\cos^{2}t] = (1/2)(1/s + s/(s^{2}+4)) = (s^{2}+2)/(s(s^{2}+4))$
  \\[8pt]\textcolor{red}{\textbf{答}}：
  $L[\sin^{2}t] = 2/(s(s^{2}+4))$\\
  $L[\cos^{2}t] = (s^{2}+2)/(s(s^{2}+4))$\\
\end{frame}

\begin{frame}{問10　解答}
  像関数の移動法則 $: L[e^{\alpha t}f(t)] = F(s-\alpha)$ より $s \to s-\alpha$ に置き換える\\[4pt]
  $L[t] = 1/s^{2}$ だから $L[te^{\alpha t}] = 1/(s-\alpha)^{2}$\\[4pt]
  $L[\cos(\beta t)] = s/(s^{2}+\beta^{2})$ だから $L[e^{\alpha t}\cos(\beta t)] = (s-\alpha)/((s-\alpha)^{2}+\beta^{2})$
  \\[8pt]\textcolor{red}{\textbf{答}}：$L[te^{\alpha t}] = 1/(s-\alpha)^{2} \ (s > \alpha)$\\
  $L[e^{\alpha t}\cos(\beta t)] = (s-\alpha)/((s-\alpha)^{2}+\beta^{2}) \ (s > \alpha)$
\end{frame}

\begin{frame}{問11　解答}
  原関数の移動法則 $: L[f(t-\mu)U(t-\mu)] = e^{-\mu s}F(s)$\\[4pt]
  ここで $f(t) = \sin t, \mu = \pi/4, F(s) = L[\sin t] = 1/(s^{2}+1)$
  \\[8pt]\textcolor{red}{\textbf{答}}：$L[\sin(t-\pi/4)\cdot U(t-\pi/4)] = e^{-\pi s/4}/(s^{2}+1)$
\end{frame}

\begin{frame}{問12　解答}
  \textbf{グラフ：}$t < 2$ では $f(t)=0$, $t \geq 2$ では原点を $t=2$ に平行移動した直線（傾き 1 ）\\[0.3em]
  $f(t) = g(t-2)U(t-2)$ と見る（ $g(t) = t$ ）\\[4pt]
  原関数の移動法則 $: L[g(t-2)U(t-2)] = e^{-2s}\cdot L[t] = e^{-2s}/s^{2}$
  \\[8pt]\textcolor{red}{\textbf{答}}：$L[f(t)] = e^{-2s}/s^{2}  (s > 0)$
\end{frame}

\end{document}
